\section{对早期数据集发布的更正}
\label{app:corrections}

早期数据集发布~\cite{survey01}以中文撰写，主要依据摘要和追踪列表中的概述。
\Cref{tab:corrections} 列出了其中为全文和在线来源所不支持的实质性陈述；AI 审查轮次中所作的其他更正随数据一并记录。

\begingroup
\small
\renewcommand{\arraystretch}{1.2}
\begin{longtable}{@{}p{0.11\linewidth}p{0.40\linewidth}p{0.43\linewidth}@{}}
\caption{对早期数据集发布的实质性更正。}\label{tab:corrections}\\
\toprule
来源 & 早期发布的表述 & 全文 / 在线来源显示 \\
\midrule
\endfirsthead
\toprule
来源 & 早期发布的表述 & 全文 / 在线来源显示 \\
\midrule
\endhead
\bottomrule
\endfoot
\cite{arx26758} & 类\jev{}开放权重模型：AUC 0.94$\to$0.23 & 此结果仅针对一个开源模型（Laya）。另一个模型（Open-Jev）有 19.5\% 的答案发生翻转。“效应随选项数量增加而增大”这一论断未得到证实。 \\
\cite{arx35342} & 恢复“unknown”后，软准确率由 0.771 升至 0.978 & 0.978 采用的是更宽松的区间指标；同口径的二值更正结果为 0.903。 \\
\cite{arx23959} & \jev{}用于诈骗筛查，零误报 & 该模型是作者基于 Qwen3-4B 的复刻版（JevLite）。在 12 个合法场景中无错误，但对模糊的合法来电，准确率为 38.1\%。 \\
\cite{arx34862} & 执行前门控的证据 & 该研究对已保存的轨迹进行回顾性分类；\jev{}在 4 个基准中的 2 个上胜出，且精确率较低。 \\
\cite{arx22664} & 决策模型能检测越界行为的证据 & 该论文不含任何类型化决策模型；其主题是测试框架（harness）保障。 \\
\cite{arx27678} & 排名在每种条件下均不同 & 在全部 12 次“条件 × 重复”请求中，按基线准确率排序与按正确性排序均不同；在托管模型中，\jev{}的基线准确率最低（77.38\%）。 \\
\cite{arx26550} & 级联以 GPT-6 费用的 41\% 提高 0.9 个百分点 & 在留出数据池上成立（RewardBench 83\%）；在 JudgeBench 上，级联为 $-0.74$ 个百分点。前瞻性实时运行以 GPT-6 费用的 75.5\% 达到与其相当的水平。所谓“盲法人工裁定”仅由一名审核者对 183 个条目完成。 \\
\cite{arx29769} & 级联仅提高 1.5 个百分点 & 在主要设置中，级联的平均提升\emph{至多}为 1.5 个百分点；在 9 个面板中有 6 个低于最佳单一评估器。 \\
\cite{arx24052} & 被列为支持校准的证据 & 原始\jev{}校准不良（斜率 1.63），经本地重新校准后才得以改善。 \\
\cite{arx24881} & 0.666 对 0.868，“依据追踪列表摘录” & 已在原文附录 E.10 中得到确认。背景：对照模型以 32{,}419 个标签训练；\jev{}的 AUROC 在数值上高于单次推理基线，而后者是在更大的数据集上评分的。 \\
\cite{arx27607} & “一致性”0.573 / 0.398 & 这些是与专家错误计数的 Kendall $\tau_b$ 相关系数；在具有临床意义的错误上，一个本地 27B 模型表现更好。 \\
\cite{arx27535} & 效应误差降低 41\% & 这是次要终点。预注册的主要终点未显示明确改善，且在第二个数据集上原始误差有所恶化。 \\
\cite{arx28919} & 基于一家拥有 10{,}000 个席位企业的真实会话模拟 & 任务是作者编写的 46 段对话；仅用户行为参数来自真实会话。\jev{}标签为实时获取，但网关并未构建；所有路由判定均在模拟器中运行。 \\
\cite{checkpoint2026jev} & 对照成本由 0.54 美元升至 4.39 美元 & 数字应为 0.56 美元升至 4.39 美元（0.54 美元是\jev{}自身的成本）。页面标题已更改。 \\
\cite{simmons2026saidno} & 在美貌、领导力、礼貌、道德四项上“白人”占 51\%；“是”回答占 3\% & 51\% 是全部八项特质的总体结果（8 项中胜出 6 项）；3\% 是“是”的平均概率。标题和日期已更改。 \\
\cite{qin2026zhdecisionbench} & 在中文中系统性过度自信；28\% 的顺序翻转 & 该基准的首个版本（9 月 26 日）未测试\jev{}（测试的是两个 Laya 模型和 Qwen3.5-2B）；28\% 来自一个 25 题的子集。其第二个版本（9 月 28 日）称过度自信这一发现部分属于小样本假象。 \\
\cite{yurin2026confident} & 置信度公式基于 76{,}807 个回答 & 该公式基于 738{,}164 个 Choice 回答得出；76{,}807 是二选项子集。 \\
\cite{typesafe_models} & 模型文档中给出 70--500\,ms & 该数字来自发布博文~\cite{almeida2026introducing}；文档所述约为 100--150\,ms。 \\
\pol{}文本 & 条款 \clause{3.4.1}、\clause{5.3.x}、\clause{3.6} & 在提交 \texttt{5791ae3} 之前，已于提交 \texttt{7750ceb}--\texttt{526da3d} 中重新编号为 \clause{5.4.1}、\clause{4.3.x}、\clause{5.6}（\cref{tab:concordance}）。 \\
\end{longtable}
\endgroup

\subsection{审查期间修订的确定性评级}
\label{sec:certainty-detail}

\cref{tab:keyfindings} 中的确定性评级最初于 10 月 4 日给出，当时只计入被引用的研究，且所用规则无法下调任何评级。
在第 12--15 轮审查期间，即更新检索的研究完成编码之后，该规则根据审查者的发现作了修订（原规则只能向上调整，未界定何为相反研究，只计入被引用的研究，也未说明一项同时研究托管模型与开源模型的研究如何计数），并以修订后的规则（\cref{sec:appraisal}）基于全部已编码的主窗口研究与更新检索研究，对每项发现重新评级。
更新检索的方案已冻结主窗口的评级，因此这构成对该方案的偏离。
最近一次修订按系统组计数证据，并且对于关于托管 \jev{} 或关于类型化模型总体的发现，仅依据托管模型证据评级。
“10 月 4 日”一栏列出 10 月 4 日给定的评级；若以现行规则对当日所引研究重新计算，其中一项将为极低（\cref{tab:certchanges} 的注~$a$）。
\Cref{tab:certchanges} 列出现行评级与 10 月 4 日给定的评级不同、或现已有争议的发现；其余评级未变。
在某项发现所涉要求上编码的每项研究，均按系统组被归为支持、相反、佐证或不在范围内，并在 \nolinkurl{data/certainty_evidence.csv} 中附有理由；\nolinkurl{scripts/rate_certainty.py} 据此计算全部三栏，并以其输出核对两张表（\texttt{--check}），\texttt{--verbose} 则打印每项发现及每个分句所计入的研究。以下各段给出非机械部分的推理。

\begin{table}[!htbp]
\centering
\small
\caption{10 月 4 日与现在之间存在差异的确定性评级（\cref{tab:keyfindings}）；“仅主窗口”为仅对主窗口研究应用现行规则的结果。\emph{10 月 4 日}：10 月 4 日给定的评级。$^{a}$~若以现行规则对当日所引研究评级则为极低：中等评级依据的是两项开源模型研究，其中一项是当日所引唯一的预注册研究~\cite{arx33843,zen23041163}，以及一个在一个数据集上达到目标、在另一个数据集上未达到目标的留出阈值~\cite{arx00346}，后者不计入任何一方；仅余一项支持性的托管研究~\cite{arx33401}。}
\label{tab:certchanges}
\begin{tabular}{@{}p{0.42\linewidth}lll@{}}
\toprule
发现 & 10 月 4 日 & 仅主窗口 & 现在 \\
\midrule
固定阈值不可迁移；需本地拟合 & 低 & 低，有争议 & 低，有争议 \\
错误集中于被总体指标掩盖的子类型 & 极低 & 低 & 低 \\
指令注入很少选中攻击者的目标 & 极低 & 极低，有争议 & 极低，有争议 \\
升级至更强的推理模型可节省成本且准确性无损失 & 低 & 低 & 低，有争议 \\
预先设定的阈值在新数据上未达到错误目标 & 中等$^{a}$ & 极低，有争议 & 低，有争议 \\
\bottomrule
\end{tabular}
\end{table}

\paragraph{更新检索研究的影响。}
与仅依据主窗口研究相比，更新检索的研究改变了一项评级：它们将“预先设定的阈值在新数据上未达到错误目标”这一发现由极低提升为低，仍有争议（见下文）。它们还使一项发现变为有争议，并为主窗口中已有争议的两项发现增加了相反研究：为注入相关发现增加一项，为关于预先设定阈值的发现增加两项；没有任何评级下降。

\paragraph{升级。}
“升级至更强的推理模型可节省成本且准确性无损失”这一发现，依据主窗口研究为低，与 10 月 4 日相同：两项针对托管 \jev{} 的较强设计研究支持该发现~\cite{arx26550,arx24574}，其中一项经预注册，另有一项中等设计研究和一项较弱设计研究与之一致~\cite{arx26532,arx02046}；第三项经预注册的较强设计研究~\cite{zen22922769}\zmark{}测试的是一个开源类型化模型，因此它佐证该发现，但不能将其提升至中等。
纳入更新检索后，该发现为低且有争议：一项较强设计研究发现，一旦计入 \jev{} 预筛查自身消耗的 token，它仅使一个 LLM 安全审核器的成本降低 4.3\%~\cite{arx02267}；该研究与该发现相反，但其排序低于经预注册的支持研究，因此评级不下降。
两项开源模型研究指向同一方向，仅作为佐证报告：一个经预注册的置信度门控级联，从一个接近随机水平的开源类型化评估器升级至 Qwen3-32B，升级了 92--100\% 的条目，因而几乎没有节省任何成本~\cite{arx07177}；将 Laya 最不确定的五分之一条目转交一个前沿模型后，在三个数据集中的两个上，macro-F1 仍低于该前沿模型自身的水平~\cite{arx03324}。
在该预注册级联中，按照 \cref{sec:appraisal} 的标准，Qwen3-32B 是更强的推理模型，其准确性（0.62--0.87）远高于那个接近随机水平的类型化评估器。
在一项重排序研究中，将 \jev{} 最不确定的配对升级至同一模型从未带来改善~\cite{arx09188}；在该研究中，Qwen3-32B 的准确性与 \jev{} 大致相当：它与人工评分者一致的频率略高（500 个配对中分别为 87.2\% 与 86.2\%，相差五个配对，该研究未报告相应检验），但以其标签训练的重排序器略差（NDCG@10 为 0.630 对 0.637；差值 0.008，95\% 区间 $-0.007$ 至 $0.022$）。
两项差异均未超出抽样误差，且方向相反，因此 Qwen3-32B 在该研究中算作同级，该研究与本发现无关。

\paragraph{两项阈值发现。}
有两项发现涉及阈值：固定阈值在不同设置之间不可迁移，因而需要本地拟合；以及预先设定的阈值在新数据上未达到错误目标。凡证据同时涉及两者之处，二者依据相同证据评级，现均为低且有争议。
一项经预注册的标注研究在托管 \jev{} 上支持这两项发现：预先固定的 0.9 置信度阈值本应使高置信条目的准确率至少达到 0.85；任务中位数为 0.815（CI 0.691--0.897），14 项任务中有 8 项未达标~\cite{arx24574}。
经预注册的网络控制研究~\cite{arx33689}在其托管组上与这两项发现相反：在 0.05 的目标下，一个经微调的开源模型的门控在 7 个变更后的问题模板中有 6 个超出目标，而 \jev{} 的门控预先设定并沿用至变更后的模板，在 7 个中有 6 个保持在目标之内，第七个达到 0.143，仍在该研究的偏移界限之内，因此 \jev{} 的预注册计数为 6 个安全、1 个在界限之内、0 个违反。一个在变更后的设置中依然成立的门控，既与“固定阈值不可迁移”的论断相悖，也同样与“阈值在新数据上未达到目标”的论断相悖；至于对一种新问题类型的动作进行认证至少需要该类型的 19 个标签，这是对认证的限制，而非迁移失败或未达目标。
对于预先设定的阈值，除该预注册研究外，另有三项针对托管 \jev{} 的较强设计研究支持该发现，其中一项来自主窗口~\cite{arx33401}，两项来自更新检索~\cite{arx03387,arx02267}；一项中等设计的更新检索研究与之一致~\cite{zen23050542}\zmark{}。除网络控制研究外，还有两项针对托管 \jev{} 的研究在其主要结果上与该发现相反：分割共形预测达到了其边际目标~\cite{zen22935043}\zmark{}；在开发集标签上设定的阈值，在此前已被观察过且带有混合标签的测试条目上，将测试假阳性率保持在 3.24\%，处于 5\% 的容差之内~\cite{zen23075657}\zmark{}。一个在一个数据集上达到其 5\% 目标（0.049）、在另一个数据集上未达到（0.058）的留出阈值不计入任何一方~\cite{arx00346}。连同该预注册研究，评级为中等；由于相反的托管组同样经过预注册且为较强设计，与最强的支持研究同级，评级因此下降一级，至低。仅依据主窗口研究，有两项研究支持该发现~\cite{arx33401,arx24574}，评级为低，而同一相反的托管组将其降至极低。
对于固定阈值，除该预注册研究外，主窗口中支持该发现的托管 \jev{} 研究为较强设计~\cite{arx29429,arx00346,arx01006,arx26550,arx39496,zen22952571}\zmark{}，另有一项中等设计和两项较弱设计的研究~\cite{arx27607,arx02046,jkf87_2026replication}\gmark{}；更新检索又增加了三项通过托管组计入或针对托管 \jev{} 的较强设计研究~\cite{arx03387,arx02267,zen22935043}\zmark{}，以及两项中等设计的研究~\cite{arx04985,zen23050542}\zmark{}。这在两个窗口中均得出中等，而相反的托管组将其降至低，有争议。
网络控制研究的开源组以及三项开源模型研究~\cite{arx33843,zen23041163,arx07177}\zmark{}支持第二项发现，但它们仅起佐证作用。
10 月 4 日，第二项发现被给定为中等。当日为其引用的研究包括：一项针对托管 \jev{} 的研究~\cite{arx33401}；一个在 CLINC 上达到其 5\% 目标（实现风险 0.049）但其界限超出该目标的留出阈值~\cite{arx00346}；以及两项开源模型研究~\cite{arx33843,zen23041163}\zmark{}，其中一项经预注册；网络控制研究在正文中被引用，但未被引作该发现的依据。原规则将这四项全部计入，得出三项较强设计研究，其中一项经预注册。按照现行规则，开源模型研究仅起佐证作用，而该留出阈值在一个数据集上达到目标、在另一个上未达到，不计入任何一方；仅余一项支持性的托管研究，评级为极低。
一个预先固定、在留出名称上依然成立的门控~\cite{zen23088660}\zmark{}既未设定错误目标，也未测试迁移，因此与两项发现均无关。

\paragraph{其他发现。}
“指令注入在托管 \jev{} 中很少选中攻击者的目标”这一发现，在主窗口中即已为极低且有争议，因为针对托管 \jev{} 的其他攻击成功率更高~\cite{arx31142,checkpoint2026jev}；更新检索增加了一项中等设计的相反研究~\cite{arx04985}，它无法再进一步下调该评级。
“类型化模型对有害内容排序良好”这一发现仍为低：支持该排序分句的研究均未经预注册，而对 \emph{Just Ask Jev}~\cite{arx29429} 的再现以较弱设计在其 44 个基准中的 5 个上重跑了作者自己的代码和数据~\cite{jkf87_2026replication}\gmark{}，因此按照 \cref{sec:appraisal}，它不是独立的重复研究；支持该发现的更新检索研究中，有两项与一项主窗口研究共享作者，并与之合计为一项研究~\cite{zen22885291,zen23075657}\zmark{}。
“插入的内容可操纵决策”这一发现为低，与 10 月 4 日相同：三项相互独立、针对托管 \jev{} 的较强设计研究结果一致~\cite{arx31142,arx01834,arx01006}，但均未经预注册或重复研究；同一作者的两项经预注册的开源模型研究结果一致~\cite{zen22959854,zen23038928}\zmark{}，但该发现涉及的是类型化模型总体，因此它们仅起佐证作用，不能提升其评级。若将其计入，评级将为中等。
“错误集中于子类型”这一发现新增了四项针对托管 \jev{} 的较强设计研究，其中两项通过其托管组计入~\cite{arx07953,arx03324,arx08675,arx10321}；“选项名称可能压过定义”这一发现新增了其中一项研究的托管组~\cite{arx07953}以及一项开源模型研究~\cite{arx02586}；但这些研究均未经预注册或重复研究，因此两项发现均仍为低。
一项经预注册的研究发现，\jev{} 在 Choice 中的偏向由选项词本身而非其键名或位置所承载~\cite{zen23179064}\zmark{}，但它与选项名称这一发现无关：这种偏向是在没有证据时朝向某个词的先验，而该词的描述与之一致，因此并没有定义被压过；并且在这一零证据框架未能通过预先声明的终止标准后，其作者已将其降级。
就校准而言，在没有检索到先例的冒犯性语言任务上~\cite{arx07953}、在 \jev{} 自身生成的游戏状态中~\cite{zen23221456}\zmark{}以及在智能体路由上~\cite{arx02267}的校准不良概率，均归入“校准是局部的”，正如主窗口中的共情任务那样~\cite{arx24574}；在 15 个车祸叙述因素中有 14 个未通过留出无偏性检验的原始概率~\cite{arx10321}亦然。
有一项研究确实发现了在熟悉任务上的校准不良：\jev{} 在以选择形式提问的情感分析任务上校准良好（ECE 0.020），但在是/否形式下（0.091）或在 AG News 主题分类上（0.089）则不然~\cite{zen22935043}\zmark{}；其主要结果是校准取决于问题形式与任务，因此我们将其归入“校准是局部的”，而非视为与第一个分句相反，第一个分句仍为低。
就“常不及最佳 LLM 准确”而言，新的比较并未指向相反方向：\jev{} 在四个仇恨言论数据集中的三个上低于最佳生成式评估器，其中一个差异显著~\cite{arx03324}，并比九种配置中的最佳者低一个百分点~\cite{arx03935}；差距小于主窗口中的差距，但并未逆转。
就“不知道”与“以上皆非”（none）而言，最有力的新研究~\cite{arx03387}（编码为 Q）发现，在答案缺失时，显式的 NONE 选项分别有 64\% 和 82\% 被选中，但随着候选项增加，被选中的频率下降；一项较弱设计研究发现，\jev{} 在 48 个没有候选项的案例中有 45 个请求新的候选项，但在 120 个不可行案例中仅升级了 4 个~\cite{arx06425}；这使托管 \jev{} 在各研究间的范围扩大到 3--94\%，符合“选择不一致”，且没有任何针对托管 \jev{} 的新研究发现“不知道”被可靠地选中。
就充分性问题而言，一项中等设计的新研究发现，托管 \jev{} 在证据充分性上的准确率为 79.8\%~\cite{arx05107}，这仅涉及第一个分句；此类问题筛选答案的效果是否不如置信度并未再次检验，因此该发现仍为极低。
就以与独立规则一致为接受条件而言，一项较强设计研究仅在词语匹配一致时才以中等置信度接受 Laya 对计划步骤的落地（grounding）（在高置信度时则无需匹配），其表现与一个预算匹配的 ReAct 智能体相当（配对差值 $-3.2$ 个百分点，CI $-7.7$ 至 $+1.2$），而输入 token 减少 83.6\%~\cite{zen22904595}\zmark{}；作为开源模型研究，它佐证第一个分句但不能提升其评级，因此两个分句均依据一项针对托管 \jev{} 的研究~\cite{arx02048}，该发现仍为极低。
